\scanpage{80}
\mspage{105}
\begin{flushright}number 371\end{flushright}

The most important case is that in which $M$ is an $A$-module \uline{free of finite dimension} $n$ ; one then canonically identifies $\handwritten{\bigwedge}(M^{\handwritten{*}})$ and $(\handwritten{\bigwedge}(M))^{\handwritten{*}\mathrm{gr}}$ by means of the canonical isomorphism $\handwritten{\lambda}_M$ ; a multivector (resp. $p$-vector) on $M^{\handwritten{*}}$ is thus identified with a linear form on $\handwritten{\bigwedge}(M)$ (resp. $\overset{p}{\handwritten{\bigwedge}}M$) and called a \uline{multiform} (resp. a $p$-\uline{form}). \handwritten{*}With this identification, formula (22) is written
\[
x_1\handwritten{\wedge}x_2\handwritten{\wedge}\ldots\handwritten{\wedge}x_n\ ,\ x_1^{\handwritten{*}}\handwritten{\wedge}x_2^{\handwritten{*}}\handwritten{\wedge}\ldots\handwritten{\wedge}x_n^{\handwritten{*}}=\det(\handwritten{\langle}x_{1}\ ,\ x_j^{\handwritten{*}}\handwritten{\rangle}).\ \handwritten{*}
\tag*{(27)}
\]

If $(e_{\unsure{i|1}})_{1\handwritten{\leq}i\handwritten{\leq}n}$ is a basis of $M$, $(e_i^{\handwritten{*}})$ the dual basis in $M^{\handwritten{*}}$, the elements
\[
e_H^{\handwritten{*}}=e_{i_1}^{\handwritten{*}}\handwritten{\wedge}e_{i_2}^{\handwritten{*}}\handwritten{\wedge}\ldots\handwritten{\wedge}e_{i_p}^{\handwritten{*}}
\tag*{(28)}
\]
(with the notation of proposition 3) form a basis of $\handwritten{\bigwedge}(M^{\handwritten{*}})$ as $H$ runs through the subsets of $\unsure{\handwritten{[}1,n\handwritten{]}|\handwritten{[\![}1,n\handwritten{]\!]}}$ ; this basis is none other than the \uline{dual basis} of the basis $(e_H)$ of $\handwritten{\bigwedge}(M)$, for it follows from (27) that if $H$, $K$ are two distinct $p$-element subsets of $\handwritten{[}1,n\handwritten{]}$, $(i_h)_{1\handwritten{\leq}h\handwritten{\leq}p}$, $(j_h)_{1\handwritten{\leq}h\handwritten{\leq}p}$ the strictly increasing sequences of the elements of $H$ and of $K$ respectively, one has
\[
\handwritten{\langle}e_{i_1}\handwritten{\wedge}e_{i_2}\handwritten{\wedge}\ldots\handwritten{\wedge}e_{i_p}\ ,\ e_{j_1}^{\handwritten{*}}\handwritten{\wedge}e_{j_2}^{\handwritten{*}}\handwritten{\wedge}\ldots\handwritten{\wedge}e_{j_p}^{\handwritten{*}}\handwritten{\rangle}=0
\]
since there exists an index $k$ \unsure{such|tul} that $j_k\handwritten{\notin}H$, hence $\handwritten{\langle}e_{i_h}\ ,\ e_{j_k}^{\handwritten{*}}\handwritten{\rangle}=0$ for every $h$.

If $M$ and $N$ are two free $A$-modules of finite type, $u:M\handwritten{\to}N$ an $A$-linear map, ${}^{t}(\handwritten{\bigwedge}(u))$ is canonically identified with $\handwritten{\bigwedge}({}^{t}u)$ by virtue of the commutativity of diagram (25) ; it follows that if $u$ is \uline{bijective}, the contragredient ${}^{t}(\handwritten{\bigwedge}(u))^{-1}$ is identified with $\handwritten{\bigwedge}({}^{t}u^{-1})$.

\uline{Remark}.- A) \uline{Extra material which must go back to} \S3 (\uline{Graded algebras}), \uline{where it was inadvertently omitted}.

Let $A$ be a commutative graded ring of type $Z$ with positive degrees, $E$ a graded $A$-algebra of type $Z$ with positive degrees, $(A_n)$ and $(E_n)$ the respective gradings of $A$ and of $E$. As an additive group, $E$ is a subgroup of the \uline{product}
% archive sheet number: 79 (preceded by a crossed-out unreadable archive number)
